\newpage
\section{Data, Information, and Knowledge}
\label{sec:data-information-knowledge}

The words \emph{data}, \emph{information}, and \emph{knowledge} are used
interchangeably in casual conversation, but they describe three distinct things.
The distinction matters in practice: confusing them leads to building systems
that produce numbers when people need answers. This section examines the
relationship between the three concepts, then argues that computing systems
operate exclusively at the data level and that the transition from data to
information is something engineers design deliberately.

\subsection{The Hierarchy}

The relationship between data, information, and knowledge is standardly
described as a hierarchy~\autocite{ackoff1989,rowley2007}. The first two
transitions are the ones this paper depends on.

Return to the tree cross-section from Section~\ref{sec:what-is-data}. A raw
count of growth rings is data: fifty-three concentric lines, nothing more. Data
are \emph{syntactic}: they have structure, but not yet meaning. Noticing that
fifty-three rings means the tree is roughly fifty-three years old, and that the
unusually narrow rings around year thirty correspond to a known regional drought,
turns that count into \emph{information}: data organized and placed in context,
in keeping with Ratzan's definition~\autocite{ratzan2004}. Information answers
questions like ``what happened?'' and ``how much?'' It is still factual and
objective, but it is now interpretable by a human being. Knowing how to read ring
width as a record of annual rainfall, reliably enough to apply the same method
to a tree no one has studied before, is \emph{knowledge}: information understood
thoroughly enough to be put to new use.

\begin{example}
\label{ex:fever}
The progression from data to knowledge can be illustrated concisely:
\begin{itemize}
  \item \textbf{Data:} \texttt{39.5}
  \item \textbf{Information:} Patient temperature is 39.5\textdegree C, recorded
        at 14:32 on 2026-06-28.
  \item \textbf{Knowledge:} This patient has a significant fever. Combined with
        readings from the past six hours showing a rising trend, the attending
        physician should be notified immediately.
\end{itemize}
The number \texttt{39.5} is a datum drawn from $\mathbb{R}$. The interpretive
layer that attaches unit, timestamp, and patient identity constitutes
information. The clinical judgment that combines this information with prior
cases and medical guidelines constitutes knowledge.
\end{example}

\subsection{What Computing Systems Do}

Computers operate almost entirely at the level of \emph{data} as defined in
Section~\ref{sec:formal-definition}. A CPU does not understand information; it
manipulates bits. A memory array does not know what its contents mean; it stores
and retrieves bytes at addresses. Shannon's foundational account of communication
is explicit on this point: the engineering problem is the reliable transmission
of \emph{symbols}, regardless of meaning~\autocite{shannon1948}. Meaning is the
problem of the sender and receiver, not the channel.

The transition from data to information is something that software engineers
design into their systems deliberately, through type systems, schemas, protocols,
and documentation. When a programmer writes \texttt{int temperature = 39;}, they
are asserting that this particular sequence of bits represents a temperature,
measured in some unit, with some precision. The meaning is in the programmer's
head and in the documentation, not in the machine. The machine stores
\texttt{0x00000027} and nothing else.

This is not a criticism of computers. It is a description of how they work, and
understanding it is the first step toward understanding why encoding schemes,
type systems, and representation formats exist: each one is an answer to the
question of how to encode information, which has meaning, as data, which does
not, in a way that preserves the meaning reliably.

\subsection{The Gap Easy Access Does Not Close}

Access to information does not supply the understanding required to know whether
that information is relevant, reliable, or correctly applied. A search engine can
return information in milliseconds; it cannot, by itself, supply the judgment
required to evaluate it. The gap between having information and having knowledge
is not closed by retrieval speed. This is an old observation, but it is worth
restating in a technical context: systems that expose data as if it were
information, or information as if it were knowledge, mislead their users in
predictable ways.

\subsection{Relationship to the DIKW Hierarchy}

Ackoff's data--information--knowledge--wisdom hierarchy~\autocite{ackoff1989}
and its many descendants in the information-science
literature~\autocite{rowley2007,zins2007} consistently treat the data tier as
their foundation without defining it formally. The function-theoretic definition
in Section~\ref{sec:formal-definition} is intended as a precise characterization
of exactly that tier, compatible with and presupposed by every DIKW-style
account. It does not compete with the hierarchy; it fills the structural gap the
hierarchy leaves open at the bottom.

Similarly, Floridi's General Definition of Information (GDI)~\autocite{floridi2010}
characterizes information as data that are well-formed, meaningful, and (in some
variants) truthful, treating data themselves as a primitive that GDI deliberately
leaves unspecified. The function-theoretic definition supplies the structural
detail that GDI leaves open, without disturbing the relationship between data and
information that GDI describes.